\begin{abstract}
We characterize finite-resource precision for estimating first-best treatment welfare in a randomized binary trial with public uniform stratum probabilities, fair assignment, consistency, and binary potential outcomes whose arm-specific means lie in \([1/4,3/4]\). Independent participants each release one message satisfying pure local differential privacy for the complete observed record. Conditional on the classical Bernstein limit for polynomial approximation of absolute value, we establish matching bounds, up to universal constants, for minimax squared error and worst-law expected length of connected confidence intervals with coverage at least \(0.90\) at every law in this causal class, including exact treatment ties. For participant count \(n\ge2\), stratum count \(d\ge2\), and privacy budget \(0<\varepsilon\le1\), let \(t=n\varepsilon^2\) be the effective private resource and \(L=\log(ed)\) the logarithmic stratum-count factor. Minimax squared error has order \(d^2/(tL)\) when \(t\ge d^2L\) and \(\min\{1,[\log(e+d^2L/t)/L]^2\}\) otherwise; honest length has the square-root order. An explicit noninteractive polynomial procedure attains both orders, and matching converses cover one-message sequential interaction with independent public randomness. The bounds characterize saturation, consistency, and the bounded-dimension boundary for inverse-resource squared error. A signed-experiment correspondence yields the same complete characterization for uniform-reference total-variation distance on a paired family, together with estimation and honest-length lower bounds on the full probability simplex.
\end{abstract}

\section{Introduction}

The population welfare attainable by assigning the better treatment in each stratum is a natural benchmark for a randomized trial. Reporting this benchmark requires estimating the larger of two conditional outcome means in every stratum and then averaging over the population. Treatment ties make this scalar target nonsmooth: a small change in a treatment contrast can change which arm supplies the maximum. When each participant reports through a confidential local channel, precision depends jointly on the number of strata, the participant count, the privacy budget, and the permitted interaction during collection.

This paper characterizes that dependence for a specified categorical trial model. Strata have public uniform probabilities, treatment is assigned fairly conditional on the stratum and both potential outcomes, observed outcomes satisfy consistency, and binary potential-outcome means belong to \([1/4,3/4]\) in every arm and stratum. Participant records are independent and identically distributed. Pure local differential privacy protects the complete observed record, including the realized stratum, assignment, and outcome. Each participant sends one message. Collection may use independent public randomness, participant-specific kernels, and standard-Borel message spaces; sequential collection additionally permits a participant's message rule to depend on earlier releases.

Under these conditions, and conditional on the classical Bernstein absolute-approximation limit, we establish matching finite-resource bounds for two reporting criteria: minimax squared-error risk and minimax worst-law expected length of connected closed intervals with coverage at least \(0.90\) at every law in the causal class. The bounds are uniform over arbitrary admissible treatment contrasts, including exact ties. An explicit noninteractive procedure attains both precision orders simultaneously, while the lower bounds apply to every admissible one-message sequentially interactive protocol.

The welfare target admits a useful decomposition. Write \(d\ge2\) for the stratum count. Let \(V(P)\), the first-best population welfare under a causal law \(P\), be the average of the larger arm mean across strata. Let \(B(P)\), the baseline observed-outcome mean under fair assignment, and \(\tau_j(P)\), the treatment contrast in stratum \(j\), denote its two components. Then \cref{obj:prop:signed-causal-bridge} gives
\[
V(P)=B(P)+\frac{1}{2d}\sum_{j=1}^d|\tau_j(P)|.
\]
The baseline is a bounded-mean target. The average absolute contrast captures the gain from choosing the better arm separately in each stratum and supplies the nonsmooth component governing the dimension-dependent precision bounds.

To state the characterization, write \(n\ge2\) for the participant count and \(0<\varepsilon\le1\) for the pure local privacy budget. Define \(t=n\varepsilon^2\), the effective private resource, and \(L=\log(ed)\), the logarithmic dimension factor. The common squared-error scale is
\[
\rho(t,d)=
\begin{cases}
\dfrac{d^2}{tL},&t\ge d^2L,\\[5pt]
\displaystyle\min\left\{1,
\left[\dfrac{\log(e+d^2L/t)}{L}\right]^2\right\},
&0<t<d^2L.
\end{cases}
\]
For both noninteractive and sequentially interactive collection, minimax squared error is comparable to \(\rho(t,d)\), and globally honest expected interval length is comparable to \(\sqrt{\rho(t,d)}\), with universal comparison constants \cref{obj:thm:uniform-private-value-frontiers}. Thus sequential adaptation improves either minimax criterion by at most a universal multiplicative factor within the declared protocol classes.

The evaluated rate gives precise resource implications. Its saturated range ends at
\[
t=\frac{d^2\log(ed)}{ed-e}.
\]
Along any admissible sequence of resource tuples, estimation error vanishes and honest intervals shrink exactly when effective private resource diverges and
\[
\frac{\log(1+d^2/t)}{\log(ed)}\longrightarrow0.
\]
Along sequences with diverging effective private resource, inverse-resource squared-error order and inverse-square-root-resource honest-length order are equivalent to eventual boundedness of dimension. When dimension diverges and \(t/d^2\) stays in a fixed positive finite interval, squared error instead has order \([\log\log(ed)/\log(ed)]^2\), with the corresponding unsquared order for honest length. The two-stratum result supplies a direct calibration across the entire resource range \cref{obj:thm:two-cell-calibration}.

The construction combines separate private estimates of the baseline and the absolute-contrast component. A sign-vector message law supplies scaled coordinate observations whose means equal the treatment contrasts. Products over distinct participants give unbiased estimates of contrast powers, and elementary-symmetric recursions compute those moments. Resource-dependent polynomial approximation estimates absolute value; in the dense branch, an independent pilot selects polynomial estimation near zero and a signed evaluation mean for larger contrasts \cref{obj:def:frontier-handle}. The analysis accounts explicitly for dependence among coordinates released in the same messages \cref{obj:lem:private-moment-and-dependence}. For the converse, moment-matched coordinate priors create separated welfare values, while higher-order transcript contraction controls their distinguishability after history-dependent privatization \cref{obj:lem:adaptive-moment-contraction,obj:lem:separated-value-mixtures}. The common approximation input is the Bernstein limit reported by \citet[Section~3.1, p.~8]{CaiLow2011Nonsmooth}.

These results complement statistical treatment choice and policy learning, which evaluate treatment rules through welfare and regret \citep{Manski2004,Kitagawa2018,Athey2021}. Here precision is evaluated for the scalar population optimum under the declared randomized-trial model. Optimal-value inference supplies closely related estimand comparisons: \citet{LuedtkeVanDerLaan2016OptimalValue} develop inference accommodating possibly nonunique optimal strategies under their conditions, and \citet[Theorem~3.3]{WhitehouseChenAusternSyrgkanis2026Softmax} establish asymptotic inference through controlled smoothing. Our characterization evaluates finite-resource minimax precision and global connected-interval honesty over the categorical causal class.

Methodologically, the paper implements the polynomial-approximation and moment-matching approach of \citet[Theorem~1 and Sections~3--5]{CaiLow2011Nonsmooth} through private participant messages. It also complements local-private functional estimation for positive power sums \citep[Theorems~2.4--2.7]{ButuceaIssartel2021Functionals} and quadratic density functionals \citep[Theorems~3.2--3.3 and~4.1--4.2]{ButuceaRohdeSteinberger2023Quadratic}. The resulting interaction comparison concerns the average magnitude of signed treatment contrasts, with its kink at equality of the arm means.

A further consequence isolates the statistical component of the welfare problem. The signed record induces a paired law with equal total mass in each pair, and its total-variation distance from the public uniform reference equals one half of the average absolute contrast. On a symmetric causal family, the protocol correspondence preserves both collection classes and the joint seed-transcript law \cref{obj:prop:signed-causal-bridge}. Conditional on the same cited Bernstein absolute-approximation limit, we obtain matching estimation and honest-length bounds on the paired family \cref{obj:thm:paired-uniform-tv-frontier}; embedding that family in the full simplex yields the corresponding lower bounds for arbitrary laws on the same alphabet \cref{obj:thm:full-simplex-tv-converse}.



The paper proceeds from related work to the trial model, privacy requirements, and reporting criteria. It then presents the welfare construction and uniform precision bounds, followed by the total-variation consequences, welfare interpretation, and limitations and future work. The appendices develop approximation and private-moment bounds, sequential contraction and decision lower bounds, a guide to the main proofs and verification scope, and proofs of the main results.
